\documentclass[a4paper,12pt]{article}
\usepackage[utf8]{inputenc}
\usepackage[spanish]{babel}
\usepackage{geometry}
\usepackage[hidelinks]{hyperref}
\usepackage{xcolor}
\usepackage{titlesec}
\usepackage{graphicx}

% Configuración de márgenes
\geometry{top=3cm, bottom=3cm, left=2.5cm, right=2.5cm}

% Estilo de títulos
\titleformat{\section}{\color{green!40!black}\normalfont\Large\bfseries}{\thesection}{1em}{}
\titleformat{\subsection}{\color{green!60!black}\normalfont\large\bfseries}{\thesubsection}{1em}{}

\title{\textbf{Manual de Usuario: Agro-Inteligencia IA}}
\author{Sistema de Monitoreo de Cultivos}
\date{\today}

\begin{document}

\maketitle
\tableofcontents
\newpage

\section{Introducción}
Este manual describe las funcionalidades del Asistente de Agro-Inteligencia. El sistema utiliza visión artificial (ESP32-CAM) y modelos de lenguaje para asistir en la detección temprana de plagas, enfermedades y deficiencias nutricionales en los cultivos.

\section{Gestión de Cultivos}

\subsection{Registro de Nuevas Parcelas/Plantas}
Para comenzar el seguimiento de una nueva unidad productiva:
\begin{itemize}
    \item \textbf{Comando:} \texttt{/nuevo\_cultivo [Nombre/Variedad]}
    \item \textbf{Ejemplo:} \texttt{/nuevo\_cultivo Tomate Cherry Invernadero 1}
    \item \textbf{Nota:} El sistema asignará un ID único (ej. \texttt{CULT-001}).
\end{itemize}

\subsection{Listado de Activos}
Para ver todos los cultivos bajo monitoreo:
\begin{itemize}
    \item \textbf{Comando:} \texttt{/cultivos}
\end{itemize}

\section{Monitoreo Visual y Diagnóstico}

\subsection{Análisis Fitosanitario (Foto)}
Para capturar una imagen y solicitar un diagnóstico automático:
\begin{itemize}
    \item \textbf{Comando:} \texttt{/foto} (o \texttt{/foto flash} si hay poca luz).
    \item \textbf{Análisis:} El agente buscará signos de clorosis, necrosis, insectos o estrés hídrico.
\end{itemize}

\subsection{Modo Vigilancia}
Para dejar la cámara monitoreando automáticamente (ej. detección de intrusos o fauna):
\begin{itemize}
    \item \textbf{Comando:} \texttt{/vigilancia on}
\end{itemize}

\section{Simulación y Entrenamiento}

\subsection{Simular Plaga}
Para probar el sistema de alertas sin tener una plaga real:
\begin{itemize}
    \item \textbf{Comando:} \texttt{/simular\_plaga [ID]}
    \item \textbf{Efecto:} Genera una alerta de "Pulgón detectado" en el registro del cultivo.
\end{itemize}

\section{Investigación Agronómica}

\subsection{Consultas Técnicas}
El agente tiene acceso a conocimientos agronómicos generales:
\begin{itemize}
    \item \textbf{Comando:} \texttt{/investigar [tema]}
    \item \textbf{Ejemplo:} \texttt{/investigar ciclo de vida de la mosca blanca en pimiento}
\end{itemize}

\section{Administración}

\subsection{Estado del Sistema}
Verifique la conexión con la cámara y el servidor:
\begin{itemize}
    \item \textbf{Comando:} \texttt{/status}
\end{itemize}

\end{document}